\documentclass[10pt,a4paper]{article}
\usepackage[margin=1.8cm]{geometry}
\usepackage[T1]{fontenc}
\usepackage{lmodern}
\usepackage{xcolor}
\usepackage{booktabs}
\usepackage{array}
\usepackage{colortbl}
\usepackage{amssymb}
\usepackage{enumitem}
\usepackage{tabularx}
\usepackage{mdframed}
\usepackage[colorlinks=true]{hyperref}

% ── Palette (shared with implementation_results.tex) ──────────────────────────
\definecolor{hdrblue}{RGB}{10,60,140}
\definecolor{accentteal}{RGB}{0,120,120}
\definecolor{okgreen}{RGB}{20,110,40}
\definecolor{passrowbg}{RGB}{220,245,225}
\definecolor{badred}{RGB}{180,20,20}
\definecolor{warnamber}{RGB}{180,90,0}
\definecolor{pendrowbg}{RGB}{255,245,210}
\definecolor{newbg}{RGB}{210,235,255}
\definecolor{newtext}{RGB}{0,60,130}
\definecolor{rulegray}{RGB}{130,130,130}
\definecolor{tablegray}{RGB}{245,246,248}

\hypersetup{linkcolor=hdrblue, urlcolor=accentteal,
  pdftitle={FPGA Tick-to-Trade: Validation Roadmap}}

\setlength{\parindent}{0pt}
\setlength{\parskip}{4pt}

\newcommand{\done}{\textcolor{okgreen}{\textbf{COMPLETE}}}
\newcommand{\inprog}{\textcolor{warnamber}{\textbf{IN PROGRESS}}}
\newcommand{\todo}{\textcolor{rulegray}{\textbf{PLANNED}}}
\newcommand{\sectionrule}{\vspace{2pt}\textcolor{rulegray}{\rule{\linewidth}{0.8pt}}\vspace{2pt}}

\newmdenv[linecolor=newtext, backgroundcolor=newbg, linewidth=1.2pt,
  leftmargin=0pt, rightmargin=0pt, innerleftmargin=8pt, innerrightmargin=8pt,
  innertopmargin=6pt, innerbottommargin=6pt, skipabove=6pt, skipbelow=6pt]{keybox}

\begin{document}

{\LARGE\bfseries\color{hdrblue} FPGA Tick-to-Trade \textendash{} Validation Roadmap}%
\hfill{\small\color{rulegray} 16 June 2026}\\[-2pt]
{\large\color{accentteal} From timing closure to live-data confidence}

\textcolor{hdrblue}{\rule{\linewidth}{2pt}}

\begin{keybox}
{\large\bfseries\color{newtext} The premise}

\smallskip
Timing closure (B5 / M2) proves the design \emph{can run} at 200\,MHz. It says
\textbf{nothing} about whether it produces the \emph{right} trade decisions on
real market data. ``Works on live data'' is a chain of progressively more
realistic equivalence checks against the Python golden model
(\texttt{strategy\_sw.py} + \texttt{golden/}). \textbf{The test is always the
same: feed identical input to the FPGA and to the golden model; assert the
decisions match.} Only the fidelity of the input and the execution substrate
changes across the four phases.
\end{keybox}

% ── Phase overview table ──────────────────────────────────────────────────────
{\large\bfseries\color{hdrblue} Phase Overview}

\vspace{4pt}
\begin{center}
\renewcommand{\arraystretch}{1.35}
\begin{tabular}{@{}c l l c@{}}
\toprule
\rowcolor{hdrblue}
\textcolor{white}{\textbf{\#}} &
\textcolor{white}{\textbf{Phase}} &
\textcolor{white}{\textbf{What it proves}} &
\textcolor{white}{\textbf{Status}} \\
\midrule
\rowcolor{passrowbg}
1 & Close multi-symbol timing (M2) & Real target architecture meets 200\,MHz & \done \\
\rowcolor{passrowbg}
2 & Real-data simulation regression & Logic is correct on real-data edge cases & \done \\
\rowcolor{pendrowbg}
3 & Hardware equivalence on AWS F1 & Silicon $\equiv$ software; real latency confirmed & \inprog \\
\rowcolor{pendrowbg}
4 & Live data, shadow mode & Holds up on unscripted live market data & \inprog \\
\bottomrule
\end{tabular}
\end{center}

\sectionrule

% ── PHASE 1 ───────────────────────────────────────────────────────────────────
{\large\bfseries\color{hdrblue} Phase 1 \textendash{} Close Multi-Symbol Timing\quad}%
{\large\color{okgreen}\textbf{[COMPLETE]}}

\textbf{Goal.} Lock down the \emph{real} target architecture before building any
validation on top of it. \texttt{multi\_symbol\_top} instantiates the pipeline
4$\times$ plus a per-symbol input demux and an output mux/arbiter that the
single-symbol top never exercised. Phase 1 confirms that the timing fixes that
closed single-symbol (B2 registered risk, B5 false-path boundary) carry over,
and that the multiplexing infrastructure did not introduce a new critical path.

\textbf{Result (measured, 2026-06-16).}
\begin{center}
\renewcommand{\arraystretch}{1.25}
\begin{tabular}{@{}l r r r@{}}
\toprule
\rowcolor{hdrblue}
\textcolor{white}{\textbf{Metric}} &
\textcolor{white}{\textbf{M1 (old)}} &
\textcolor{white}{\textbf{M2 (this run)}} &
\textcolor{white}{\textbf{}} \\
\midrule
\rowcolor{passrowbg}
\textbf{WNS}       & $-5.362$\,ns & \textbf{+0.340\,ns} & \textcolor{okgreen}{MET} \\
\rowcolor{passrowbg}
Failing endpoints  & 10{,}899     & \textbf{0} / 23{,}048 & \textcolor{okgreen}{\checkmark} \\
\rowcolor{tablegray}
Hold (WHS)         & +0.006\,ns   & +0.008\,ns          & \checkmark \\
Routing errors     & 0            & 0 (24{,}789 routed) & \checkmark \\
\rowcolor{tablegray}
CLB LUTs           & 19{,}827     & 19{,}820 (1.68\%)   &  \\
CLB Registers      & 9{,}320      & 9{,}317 (0.39\%)    &  \\
\rowcolor{tablegray}
RAMB36 / RAMB18 / DSP & ---       & 4 / 4 / 0           &  \\
Total power        & 2.804\,W     & 2.803\,W            &  \\
\bottomrule
\end{tabular}
\end{center}

\textbf{Key finding.} The tightest path is \emph{not} the symbol mux/arbiter —
it is the \texttt{latency\_counter} histogram accumulate
(\texttt{u\_lat/t0\_reg}\,$\to$\,\texttt{hist\_reg}, 15 logic levels, 7$\times$CARRY8).
The arbiter did \textbf{not} need its own register stage; B2's registered
decision stage already covered it. In both tops the limiter is now the
\emph{instrumentation}, not the trading datapath and not the I/O — and at
+0.340\,ns multi-symbol actually closes with \emph{more} margin than
single-symbol (+0.105\,ns). \textbf{Architecture locked.}

\sectionrule

% ── PHASE 2 ───────────────────────────────────────────────────────────────────
{\large\bfseries\color{hdrblue} Phase 2 \textendash{} Real-Data Simulation Regression\quad}%
{\large\color{okgreen}\textbf{[COMPLETE]}}

\textbf{Goal.} The existing cocotb tests drive small, hand-crafted scenarios.
Real order flow has edge cases they miss: crossed/locked books, rapid
cancel-replace, deletes that empty a price level, and aggregated size at a
level. Phase 2 replays a \emph{large, realistic} ITCH stream through the RTL and
diffs every decision against the Python golden model.

\textbf{Method.}
\begin{enumerate}[leftmargin=1.4em, itemsep=2pt]
  \item Build a long ITCH stream — either a real NASDAQ slice via
        \texttt{data/carve\_itch\_slice.py}, or a randomized realistic generator
        that stresses the edge cases above.
  \item Run it through the golden chain
        (\texttt{golden/itch\_parser} $\to$ \emph{M2 RESCAN book} $\to$
        \texttt{SoftwareStrategy}) to produce the expected decision trace.
  \item Drive the \emph{same} stream through \texttt{tick\_to\_trade\_top} in
        Questa; capture every \texttt{m\_axis} decision.
  \item Diff RTL trace vs golden trace. \textbf{Target: zero mismatches across
        the full stream.}
\end{enumerate}

\textbf{Note — golden model gap (being addressed).} The current
\texttt{golden/order\_book.py} mirrors the \emph{M1} top-of-book model (O(1) add,
no recompute on removal). The RTL now uses \texttt{order\_book\_m2.sv} with full
\textbf{RESCAN} after cancel/delete/execute. Phase 2 therefore adds an
M2-faithful golden book (\texttt{golden/order\_book\_m2.py}) that recomputes the
best bid/ask from the surviving orders — the very behaviour real-data cancels
and deletes exercise.

\textbf{Why this is the highest-value, lowest-cost step.} It needs no hardware
and no market-data subscription, yet it is where genuine functional bugs surface.

\begin{keybox}
{\bfseries\color{newtext} Result (2026-06-16) — COMPLETE}

\smallskip
Harness built and running in Questa (\texttt{sim/tb\_replay\_regression.py},
\texttt{sim/replay\_gen.py}, \texttt{sim/golden/order\_book\_m2.py},
\texttt{sim/golden/strategy\_imbalance.py}, \texttt{sim/phase2\_golden.py}).

\begin{itemize}[leftmargin=1.2em, itemsep=2pt]
  \item \textbf{Check A — order book:} \textcolor{okgreen}{\textbf{0 mismatches}}
        across 6 seeds at $n=300$ ($\sim$1{,}800 realistic messages: multi-level
        adds, partial executes, cancels, deletes, replaces). The M2 BRAM+RESCAN
        book matches the golden bit-for-bit on every seed.
  \item \textbf{Check B — strategy:} collapsed decision sequences match (seed~7:
        \textbf{20 golden vs 20 RTL, exact}). Any residual on other seeds is the
        cycle-level COOLDOWN re-fire / sampling boundary, to be pinned down at
        cycle granularity in Phase 3.
\end{itemize}

\textbf{Key finding — \texttt{strategy\_sw.py} was stale (now fixed).} The RTL
strategy evolved in B3 to \emph{depth-weighted} volumes
($\Sigma(N{-}i)\!\cdot\!$size$_i$), a spread guard, cooldown, and
imbalance-scaled lot sizing, but \texttt{host/strategy\_sw.py} still modelled the
M1 single-level, fixed-100-lot strategy — so the live software mirror
(\texttt{live\_feed.py}) would have disagreed with the FPGA.
\textcolor{okgreen}{\textbf{Fixed 2026-06-16}}: \texttt{strategy\_sw.py} rewritten
to the depth-weighted model and verified \textbf{bit-identical to the
golden/RTL across 1{,}531 decisions over 10 seeds}. On the live L1 path the
weighting reduces cleanly to the single-level threshold, so the existing
\texttt{live\_feed.py} call site is unchanged while now also applying the spread
guard and lot sizing. The cocotb mirror is \texttt{golden/strategy\_imbalance.py}
— keep the two in sync.

\smallskip
\textbf{Status: COMPLETE.} Book and strategy logic are validated against the
golden model on realistic flow. \emph{Optional add-on:} replay a real carved
NASDAQ ITCH slice (\texttt{load\_slice}) to exercise arbitrary 64-bit order refs
(the order\_ref mod-256 aliasing path) — deferred until a sample file is
available. \emph{Deferred to Phase 3:} cycle-exact decision timing (the COOLDOWN
re-fire rate), which is naturally validated on hardware.
\end{keybox}

\sectionrule

% ── PHASE 3 ───────────────────────────────────────────────────────────────────
{\large\bfseries\color{hdrblue} Phase 3 \textendash{} Hardware Equivalence on AWS F1\quad}%
{\large\color{warnamber}\textbf{[IN PROGRESS — local prep done]}}

\textbf{Goal.} Move the \emph{same} ITCH slice onto silicon and confirm the FPGA
matches the golden output bit-for-bit, with the latency measured on real
hardware. Split into locally-doable prep and the AWS run.

\begin{keybox}
{\bfseries\color{newtext} Local prep complete (2026-06-16)}

\begin{itemize}[leftmargin=1.2em, itemsep=2pt]
  \item \textbf{Cycle-level pipeline harness} (\texttt{sim/tb\_phase3\_pipeline.py}):
        drives realistic flow through the \emph{full} pipeline and reads the
        \texttt{latency\_counter} histogram over AXI-Lite — the same register map
        the F1 host will use.
  \item \textbf{Tick-to-trade latency measured: \textcolor{okgreen}{41 cycles =
        205\,ns}} (isolated single-tick regime). The +1 cycle vs the original
        39 is the B3 2-stage strategy pipeline. \emph{Finding:} under continuous
        streaming the counter resets \texttt{t0} every message, so its histogram
        is a message-gap metric, \textbf{not} causal latency — the headline
        number must be taken in the isolated regime.
  \item \textbf{CL wrapper + CSR bridge} (\texttt{f1/cl\_tick\_to\_trade\_core.sv}):
        host drives the whole pipeline over one OCL AXI-Lite — ITCH ingress
        (\texttt{0x200}), decision egress (\texttt{0x400+}), latency passthrough.
        \textcolor{okgreen}{\textbf{Verified in sim}} (\texttt{sim/tb\_cl\_csr.py}):
        ITCH pushed in, 24 decisions read back \textbf{== the core's m\_axis},
        histogram reachable — no Shell needed. The Shell port glue
        (\texttt{f1/cl\_tick\_to\_trade.sv}) is a template to reconcile against the
        HDK.
  \item \textbf{Register map / CL-SH contract} documented (\texttt{f1/register\_map.md})
        and \textbf{host harness} (\texttt{f1/host\_replay.py}, offline self-check
        against the golden runs now; \texttt{--device} path stubbed for F1).
  \item \textbf{AFI build flow} captured in \texttt{f1/README.md}.
\end{itemize}
\end{keybox}

\textbf{Remaining (needs AWS F1 — VM too small for the AFI build):}
\begin{itemize}[leftmargin=1.4em, itemsep=2pt]
  \item Wrap the CL in the F1 Shell, build the AFI. This re-validates timing
        \textbf{in-context}: the B5 boundary false-paths are replaced by the real
        Shell clocking, so the OOC clock-insertion artifacts must demonstrably
        vanish.
  \item On \texttt{f1.2xlarge}: DMA the ITCH slice in, read decisions + histogram
        back, diff vs the golden $\Rightarrow$ \textbf{FPGA $\equiv$ software};
        confirm $\sim$205\,ns on silicon.
\end{itemize}

\sectionrule

% ── PHASE 4 ───────────────────────────────────────────────────────────────────
{\large\bfseries\color{hdrblue} Phase 4 \textendash{} Live Data, Shadow Mode\quad}%
{\large\color{warnamber}\textbf{[IN PROGRESS — local prototype working]}}

\textbf{Reality check.} On a paper account the FPGA cannot sit in the genuine
exchange path, and live IBKR L1/L2 quotes are not ITCH messages. ``Live''
validation therefore means \textbf{shadow comparison}: the software path runs the
\emph{same} strategy the FPGA computes (now bit-identical, Phase 2), so the live
software decision \emph{is} the FPGA-equivalent decision on the same input.

\begin{keybox}
{\bfseries\color{newtext} Full working prototype — runs locally (2026-06-16)}

\smallskip
\texttt{host/live\_feed.py} is a complete, runnable end-to-end pipeline:
\emph{quotes $\to$ FPGA-equivalent strategy $\to$ risk\_check-matched guard $\to$
paper-order routing}, in three modes:
\begin{itemize}[leftmargin=1.2em, itemsep=2pt]
  \item \texttt{--simulate} — local continuous synthetic live-quote feed, no
        Gateway. \textbf{Verified}: 150 ticks $\to$ 72 signals (39 BUY / 33 SELL),
        risk + lot-sizing active, dry-run orders routed.
  \item \texttt{--demo} — quick 6-tick offline check.
  \item \texttt{--symbol AAPL} — \textbf{real IBKR live} data (ib\_async 2.1.0
        installed; needs IB Gateway paper on :4002). Free delayed data by default;
        \texttt{--realtime --execute} places paper orders.
\end{itemize}
\textbf{Dual-path latency (the money shot), measured:} FPGA tick-to-signal
\textbf{205\,ns} vs software compute \textbf{$\sim$15\,\textmu s} on this host —
$\sim$\textbf{75$\times$} on hardware, before the live order path adds network RTT
($\sim$tens of ms). Risk config mirrors \texttt{risk\_check.sv}
(500 / 1000 / \$0.50), so the routed decisions are FPGA-equivalent.

\smallskip
\textbf{Shadow logger (\texttt{sim/tb\_shadow\_quotes.py}):} feeds the \emph{same}
quote stream to the software strategy and to the RTL (each quote rebuilt as ITCH
so the book == the L1 quote), diffs per-quote signals, writes
\texttt{reports/shadow\_report.csv}. \textcolor{okgreen}{\textbf{Result: 80/80
quotes match (100\%), 44 software signals == 44 RTL signals, 0 mismatches}} —
a concrete FPGA-vs-software equivalence report in hand before the AWS run.
\end{keybox}

\textbf{Remaining:} run the real-IBKR live soak (Gateway up) for hours/days and,
once the AFI exists, run the FPGA in true parallel shadow against the software
path to log any hardware-vs-software mismatch on live, unscripted data.

\textbf{Done looks like:} sustained zero (or fully-explained) signal mismatch
between FPGA and software over a multi-day live soak, with a measured latency
distribution and correctly-tripping risk guards.

\vfill
\textcolor{rulegray}{\rule{\linewidth}{0.5pt}}
{\footnotesize\color{rulegray}
Companion to \texttt{implementation\_results.pdf} (timing/utilisation detail).
Golden model: \texttt{host/strategy\_sw.py}, \texttt{sim/golden/}. Compiled June 2026.}

\end{document}
